\documentclass[11pt,a4paper]{article}
\usepackage[T1]{fontenc}
\usepackage[utf8]{inputenc}
\usepackage{lmodern,microtype}
\usepackage[margin=25mm,headheight=14pt]{geometry}
\usepackage{amsmath,amssymb,amsthm,mathtools}
\usepackage{booktabs,array,longtable,enumitem,xcolor}
\usepackage{fancyvrb,fancyhdr}
\usepackage[breaklinks,colorlinks,linkcolor=blue!45!black,urlcolor=blue!45!black,citecolor=blue!45!black]{hyperref}
\hypersetup{pdftitle={Proof compression for the ZFC busy-beaver construction},pdfauthor={Research collaboration for ProveIt},pdfsubject={FM215, exact-table Lean equivalence, decoder certificates and bounded 214-state search}}
\RecustomVerbatimEnvironment{verbatim}{Verbatim}{fontsize=\small}
\setlist{itemsep=3pt,topsep=5pt}
\setlength{\emergencystretch}{3em}
\setlength{\parskip}{4pt}
\pagestyle{fancy}
\fancyhf{}
\fancyhead[L]{\small ProveIt: ZFC machine proof compression}
\fancyhead[R]{\small 21 October 2026}
\fancyfoot[C]{\thepage}
\newcommand{\Con}{\operatorname{Con}}
\newcommand{\Halts}{\operatorname{Halts}}
\newcommand{\repo}[2]{\href{https://github.com/VladimirReshetnikov/ProveIt/blob/b6c7087477e51c017003bd729e712c8969733223/#1}{#2}}
\newcommand{\zfcroot}{Computability/BusyBeaver/Research/ZFC/}
\newtheorem{proposition}{Proposition}[section]
\title{Proof compression for the ZFC\ busy-beaver construction\\[5pt]
\large A new four-row 215-state variant, checked local equivalence,\\
context certificates, and limits of the 214-state search}
\author{Research packet for Vladimir Reshetnikov's ProveIt repository}
\date{21 October 2026}

\begin{document}
\maketitle
\begin{abstract}
The most economical way to certify nearby machines is to prove a small
whole-tape simulation against the already certified primary K215 and reuse
its calculus calibration. We construct a distinct 215-state machine,
FM215, by changing four transition rows. Two immediate incoming-edge facts
give its invariant; every parent transition is simulated by one or three
new transitions. The generic transfer theorem and the exact-table,
all-initial-tapes equivalence have been compiled in Lean 4.32.0 with only
\texttt{propext} and \texttt{Quot.sound} as reported dependencies.

For shortening the parent proof itself, we identify a remaining Rocq
certificate-symmetry port and construct a 234-role finite certificate for
the parent's 4,750 internal decoder-prefix occurrences. The latter is an
executable internal-decoder certificate, not yet a Lean or Rocq runtime
replacement. A separate compiled Lean countermodel excludes a logic-only
recovery step proposed for the large-matrix experiment, without deciding
that experiment's full equiconsistency. A bounded search of 1,761,648 local
replacement rows across 15 machines, together with a separate helper
search, yields no verified 214-state machine. The packet includes exact
tables, proofs, scripts, receipts, source pins and inherited notices.
\end{abstract}

\section*{Decisions and scope}
\addcontentsline{toc}{section}{Decisions and scope}
\begin{enumerate}
\item \textbf{Use FM215 for a compact incremental correctness proof.}
  It has a simpler invariant than the existing conditional-marker CM215.
  It remains a 215-state machine: the removed logical role is reused as
  a physical selector, so no physical state is saved.
\item \textbf{Refactor the shared proof boundary before regenerating large
  certificates.} Positive simulation, one-way certificate symmetry, and
  context-indexed decoder sharing target actual repeated obligations.
\item \textbf{Keep the 214-state question open.} No 214-state candidate
  is offered as correct. The search limits and the logical obstruction
  are explicitly restricted results, not lower bounds.
\end{enumerate}

\clearpage
\tableofcontents
\clearpage

\section{What is proved, and which machine is meant}
\label{sec:status}

All repository comparisons use the fixed commit
\begin{center}\texttt{b6c7087477e51c017003bd729e712c8969733223}.\end{center}
Here ``primary K215'' means the native table at
\path{candidate-k215/k215.native.tm}, whose exact hash appears in
Section~\ref{sec:fm-proof}. The machine convention is one two-way infinite
binary tape, initially all zero; every transition moves left or right;
the halt label is separate from the working-state count. Initialization
is part of the machine. FM215 has 430 working transition rows and 215
working states. Its native start, dispatcher and halt labels are
$0,84,254$; its dense labels are $0,\ldots,214$, with start $0$,
dispatcher $83$ and halt $215$. The supplied bijection specifies the
renaming exactly.

Let $\operatorname{Ref}_{\mathrm{ZFC}}$ denote existence of a finite
refutation in the ZFC calculus calibrated by the repository, and define
$\Con(\mathrm{ZFC})=\neg\operatorname{Ref}_{\mathrm{ZFC}}$. The inherited
primary theorem and the new ordinary simulation argument compose as
\[
 \Halts(\mathrm{FM215},\mathbf 0)
 \;\longleftrightarrow\;
 \Halts(\mathrm{K215},\mathbf 0)
 \;\longleftrightarrow\;
 \operatorname{Ref}_{\mathrm{ZFC}}.
\]
Negating this equivalence gives the requested
\[
 \mathrm{ZFC}\ \vdash\
 \neg\Halts(\mathrm{FM215},\mathbf 0)
 \;\longleftrightarrow\;\Con(\mathrm{ZFC}).
\]
The added argument consists of finite table facts and induction on runs
and is formalizable in ZFC. This statement does not assume or establish
$\Con(\mathrm{ZFC})$, and therefore does not assert unconditional
nonhalting. As with the repository's documented verification boundary,
an external proof-assistant theorem is distinguished from an explicitly
encoded object-theory derivation. This packet does not emit the code of a
ZFC or PA proof of the displayed statement.\cite{verification}

\begin{longtable}{@{}>{\raggedright\arraybackslash}p{.26\textwidth}>{\raggedright\arraybackslash}p{.35\textwidth}>{\raggedright\arraybackslash}p{.33\textwidth}@{}}
\toprule
Claim & Evidence in this packet & Exact limitation\\
\midrule
\endhead
FM215 and primary K215 equihalt from state 0 on every tape &
Compiled Lean theorem, exact literal-table binding audit, ordinary proof &
Functional tapes; the adapter to the repository's zipper/dense endpoint
is described but not compiled here.\\[5pt]
FM215 has the same calibrated blank entry &
Complete synchronized 95,714/95,840-step prefix with identical whole tape &
Supporting check; unbounded equivalence follows from the Lean simulation.\\[5pt]
Internal decoder has 234 context roles &
468 checked branch rules and source-DAG expansion equality &
Primitive first-row binding only; primitive runtime remains to be proved.\\[5pt]
Missing low-binder vacuity is not logically recoverable &
Ordinary countermodel and compiled Lean named logical interface &
Excludes nonlogical seed bundles; no full-calculus inconsistency claim.\\[5pt]
No successful 214-state reduction in tested families &
Reproducible search scripts and counts &
Bounded templates and overapproximate invariants; no machine lower bound.\\
\bottomrule
\end{longtable}

\section{Audit of the existing formal proofs}
\label{sec:proof-audit}

The inspected ProveIt revision is
\texttt{b6c7087477e51c017003bd729e712c8969733223}.
The repository already contains three proof-simplification checkpoints.
Consequently, deriving a shared AX-word interpreter, removing the reverse
Lean source certificate, or replacing the original expanded dominated-rule
derivation would repeat work already present in this revision.\cite{verification,phase3}

The current Lean primary dependency closure has 185 project modules,
compared with 204 at the original full-build checkpoint. The latest
phase-three record freshly rebuilds nine affected modules and binds the
176 reused objects to their earlier checked sources, dependencies, objects
and logs. Its two explicit audit clients produce ten reports containing
only \texttt{propext}, \texttt{Classical.choice}, and \texttt{Quot.sound}.
This is an incremental check of the complete dependency closure, not a
fresh compilation of all 185 modules in this session. The original Rocq
record reports 237 selected sources, four closed endpoint assumption
reports and a recursive \texttt{coqchk}. Subsequent simplification records
retain that baseline and identify the affected rebuilt modules.

The precise endpoint in Lean is
\begin{verbatim}
KEndpoint215.blank_halts_iff_zfc_refutation :
  BlankHalts machine215 <-> BProv ZFCax_s [] Form.fBot
\end{verbatim}
The Rocq dense endpoint is
\texttt{tm215\_halts\_iff\_ZFC\_inconsistent} in
\path{k225-formal/rocq/rocq-extension/coq/215/CorrectDense215.v}.
Both refer to the primary machine's literal dense table. The core
composition is already short: the Lean endpoint is a direct application
of the shared \texttt{KEndpoint} theorem, and Rocq rewrites the dense/sparse
equivalence and applies its sparse endpoint. Shortening those final lines
has little mathematical value. The substantial work is in the reusable
simulation and calculus layers, and in the machine-specific certificates.\cite{endpoint}

\begin{center}
\begin{tabular}{lrr}
\hline
Proof source & Original bytes & Current bytes\\
\hline
Lean dominated implication & 689,857 & 35,185\\
Lean AX routine & 202,222 & 7,563\\
Lean source data & 653,569 & 456,548\\
Lean site-coverage facade & 107,243 & 14,916\\
Lean register determinism & 17,209 & 12,988\\
Rocq dominated implication & 921,691 & 22,673\\
\hline
\end{tabular}
\end{center}

These are the repository's recorded source-byte comparisons, rather than
measurements of the size of the resulting kernel proof terms. They also
do not count every retained historical receipt as removed repository data.\cite{simplification}
The current Lean dominated-rule argument has a 61-lemma bootstrap,
81 reused conditional schemas, eight numerical dominance hypotheses, and
a deduction theorem. The 629-node/735-instruction certificate remains
the checked Rocq rendering; it is not the current Lean rendering.

The audit sources are the primary verification README, the simplification
README, the phase-three README and its Lean record, and the actual endpoint
files. Useful current source anchors are
\texttt{KEndpoint215.lean:43} (the endpoint),
\texttt{KSim.lean:317} (positive-step halting transfer),
\texttt{RMOffsetBisim.lean:287} (one certificate gives equivalence),
\texttt{L2AxiomsPerm.lean:100} (AX-word composition), and
\texttt{DomRho8.lean:597} (the unrestricted dominated-row conclusion).
All of these anchors refer to the pinned revision above.

\subsection{A concrete remaining simplification in Rocq}

The Lean source correspondence now uses a single forward certificate and
the mathematical symmetry of an offset bisimulation. The current Rocq
\texttt{Source215.v} still separately evaluates
\begin{verbatim}
certificate_forward : ocheck (ptree p215) (ptree dagPerm) Cert215 = true
certificate_backward :
  ocheck (ptree dagPerm) (ptree p215) (omirror Cert215) = true
\end{verbatim}
with \texttt{vm\_compute} (lines 8--13), and the same directory's
\texttt{OffsetBisim230.v:167}, \texttt{ocheck\_iff}, requires both results.
Thus the same mathematical
improvement remains available in the Rocq lane.

The appropriate addition is the generic theorem
\begin{verbatim}
ocheck_mirror t1 t2 C :
  ocheck t1 t2 C = true ->
  ocheck t2 t1 (omirror C) = true.
\end{verbatim}
Its proof has five small ingredients. Lookup in \texttt{omirror C} is
lookup in \texttt{C} followed by swapping the two paths and negating the
integer offset. The identity
$\delta(-z,i)=-\delta(z,i)$ turns each old count equation into its reverse.
A simple completeness lemma for \texttt{counts\_ok}, paired with the
existing \texttt{counts\_ok\_sound}, avoids reasoning about the order of
the duplicate register indices in its test list. These facts transfer a
branch check, then an entry check, then the initial and universal checks.
The initial offset is zero, so it is fixed by negation. This removes the
second whole-certificate evaluation while preserving the exact theorem
and every table/source input. This Rocq change is a proposed port; the
new files in this packet do not claim to have compiled it.

\subsection{Use a verified parent as the proof boundary}

For a small transition-table edit, the most effective next simplification
is often a direct equivalence proof with the already verified machine.
The entire numerical-calculus calibration, proof enumeration, source
bisimulation, site coverage and blank initializer can then remain behind
one proved interface. The new obligation is a local simulation theorem
about the changed transition graph.

This applies especially well to the new FM215 construction below. The
existing conditional-marker alternative CM215 also has an ordinary
all-tape equivalence with the primary machine, but its invariant needs
an extra fact about state 156 and a two-step incoming path. FM215 uses two
unconditional predecessor-marker facts, each established by the
immediately preceding transition. Its additional correctness proof is
therefore structurally simpler.

This comparison is about the additional proof after the primary endpoint
has been proved. It is not a claim that changing the table removes the
underlying set-theoretic calibration. Likewise, the rebalanced-address
alternative's maximum source depth of 41 can reduce some certificate
arithmetic, but a shorter address is not itself a shorter correctness
argument. The composed allocation shares the primary source and seed,
which makes it a conservative candidate for reusing the existing
certificate pipeline, but it still needs its own physical runtime
certificates unless an additional direct equivalence is established.

\section{A reusable positive simulation theorem}
\label{sec:positive-simulation}

Let $f:X\to X\cup\{\bot\}$ and $g:Y\to Y\cup\{\bot\}$ be deterministic
partial transition functions. A configuration is terminal precisely when
its transition is $\bot$. Let $R\subseteq X\times Y$ satisfy:
\begin{enumerate}
\item If $R(x,y)$ and $f(x)=x'$, then some $y'$ satisfies $R(x',y')$
  and $y\mathrel{\longrightarrow_g^{k}}y'$ for some $k\geq1$.
\item If $R(x,y)$ and $f(x)=\bot$, then the $g$-computation from $y$
  halts after finitely many steps.
\end{enumerate}
Then $R(x,y)$ implies that $f$ halts from $x$ if and only if $g$ halts
from $y$. The relation need not be injective, surjective, or functional.
Reachability invariants can simply be included in $R$.

For the forward implication, induct on the finite source computation,
concatenating the finitely many target blocks. For the reverse implication,
induct on the number $t$ of successful target steps before termination.
If the source is terminal, the conclusion is immediate. Otherwise its
next step is simulated by a block of length $k\geq1$. This block cannot
extend past the target's terminal configuration: determinism and the
absence of a transition there would contradict its certified successful
completion. Hence $k\leq t$, and the induction hypothesis applies at the
next related boundary with the strictly smaller bound $t-k$.

Positive progress is essential. If zero-length target blocks were
allowed, a nonhalting source self-loop could be related to a terminal
target forever. No uniform upper bound on block lengths is needed for
halting equivalence. In FM215 there is nevertheless the useful exact
bound $k\in\{1,3\}$.

The supplied \texttt{PositiveSimulation.lean} is a generic extraction of
the final argument in the repository's \texttt{KSim.halts\_iff}. It uses
arbitrary configuration types and deterministic partial transitions, so
it can also serve machine-to-machine optimizations. It is not presented
as a new abstract computability theorem. Its value is as a small reusable
formal interface. The accompanying execution receipt records a successful
Lean 4.32.0 compilation of this theorem. Its printed dependencies are
exactly \texttt{propext} and \texttt{Quot.sound}.

\section{FM215: the short correctness argument}
\label{sec:fm-proof}

Write $T$ for the primary native K215 table and $N$ for FM215. Define the
state map
\[
 \phi(127)=124,\qquad \phi(q)=q\quad(q\neq127).
\]
The two relevant old states have rows
\[
\begin{array}{c|cc}
 &0&1\\\hline
124&0R62&1R138\\
127&0R14&1R138.
\end{array}
\]
Their entire incoming-edge census is
$104/1=1R124$ and $110/0=0R127$. Hence, on every old computation from
initial control zero and on every initial tape,
\[
 q=124\Longrightarrow u(h-1)=1,\qquad
 q=127\Longrightarrow u(h-1)=0,
\]
where $u$ denotes the current tape and $h$ the current head position.
Call this two-clause assertion $\mathcal I$. It holds initially because control
zero is neither exceptional state, and it holds after \emph{every}
old transition: the two possible destination states force their marker
through their unique incoming edge. In particular its preservation
requires no hypothesis about the preceding tape.

FM215 changes exactly four rows:
\[
\begin{array}{c|cc}
\text{row}&\text{primary K215}&\text{FM215}\\\hline
124/0&0R62&0L127\\
127/0&0R14&0R8\\
127/1&1R138&1R109\\
110/0&0R127&0R124.
\end{array}
\]
The existing provider rows are $8/0=0R14$ and $109/0=0R62$.
At simulation boundaries relate configurations with the same entire tape
and head, old invariant $\mathcal I$, and new control $\phi(q)$.
An ordinary old transition is simulated by one new transition. On the
common read-one branch, both old exceptional states go to 138, and their
new representative 124 does the same.

On an exceptional read-zero branch, let the predecessor marker be $m$.
The new computation is
\[
 (124,h)\longrightarrow(127,h-1)
 \longrightarrow(x_m,h)\longrightarrow(a_m,h+1),
\]
where $(x_0,a_0)=(8,14)$ and $(x_1,a_1)=(109,62)$.
The three writes are respectively $0,m,0$, each equal to the value
already present. Thus the entire tape is unchanged. The invariant picks
the correct old destination, and the head has moved right exactly once.
All three internal controls have a transition, so no halt is hidden in
the block. The state map fixes the terminal label 254.

The positive simulation theorem now proves halting equivalence on all
binary initial tapes, with any initial head position, from control zero.
If $z(t)$ counts the old read-zero visits to states 124 or 127 in a
successful prefix of length $t$, the exact corresponding new duration is
\[
 t+2z(t),\qquad t\leq t+2z(t)\leq3t.
\]
This is a statement about arbitrary whole tapes, not a finite sample of
register values and not an inference from the blank prefix.

\subsection{How the formal certificate binds the machine}

\texttt{FM215Data.lean} embeds both complete native tables and defines
functional tapes $\mathbb Z\to\{0,1\}$. The ambient label type is
\texttt{Fin 512}, because the native tables use labels up to 302. This
is an indexing device: exactly 215 labels have working rows, and every
outgoing target is a working label or the separate terminal label 254.
The finite ambient type does not assert a 512-state machine.

\texttt{FM215.lean}'s \texttt{incoming} theorem checks the two required
incoming-row properties for every ambient state and both symbols.
\texttt{outer} checks all unchanged and redirected rows under $\phi$.
Both are finite kernel decisions. The subsequent proofs establish
invariant preservation, the complete three-step tape identity, and the
one-or-three-step simulation. The endpoint is
\begin{verbatim}
FM215.all_tapes_halts_iff (h : Int) (t : Tape) :
  Halts (step old) (0,h,t) <-> Halts (step new) (0,h,t).
\end{verbatim}
The actual file uses the corresponding structure notation and qualified
\path{PositiveSimulation.Halts}. The blank instance is obtained by
$h=0$ and the constantly false tape.

All three supplied modules were compiled successfully in this session
with the official Lean 4.32.0 release, commit
\texttt{8c9756b28d64dab099da31a4c09229a9e6a2ef35}. The generic theorem
and both FM215 endpoints report exactly \texttt{propext} and
\texttt{Quot.sound}; there is no admitted theorem, extra assumption, or
native evaluation oracle. The final compiler logs contain no warnings.
\texttt{fm215-kernel-receipt.json} binds the checked sources, compiled
objects and logs by SHA-256. The successful sequential run took about
26 seconds in this environment; this is a reproducibility datum rather
than a comparative benchmark. The generic proof file has 99 lines,
the FM215 proof file 172, and the table-and-semantics file 918, counting
comments and blank lines. The first two files together are 10,375 bytes.
\texttt{check\_fm215.py} reproduces the compilation and the binding audit
without any dependency on the ProveIt project build.

The independent \texttt{check\_lean\_bindings.py} reads every literal Lean
row and compares it with the exact \texttt{.tm} input. It also checks the
ambient bound, so a large numerical literal cannot silently wrap under
\texttt{Fin}'s numeral instance. The primary native SHA-256 is
\begin{verbatim}
9683d6ed9628e6ce407c315b08f5da3d063da94e92706f275ca0f33e8272e4b2
\end{verbatim}
and the FM215 native SHA-256 is
\begin{verbatim}
3ece1233f11d33375425221c5bcedba941e073763912990eabf321197772dce5
\end{verbatim}
The data audit is separate from the Lean compiler check. An exact-table
theorem in functional-tape semantics is also separate from composing the
repository's existing \texttt{KEndpoint215} theorem, whose interface uses
its own dense and zipper semantics. The outstanding imported bridge has
two specific parts: identify the old native table with
\texttt{KTable215.lean}'s \texttt{tm215} on working and terminal controls
under the repository's renaming $254\mapsto230$, $301\mapsto102$,
$302\mapsto120$; and relate functional tapes to \texttt{KZip}'s two
finite lists, proving the one-step and terminal cases. The existing
\texttt{KDenseBridge} already supplies the dense-to-zipper side.
Alternatively one can adapt that generic bridge directly from Core's
finite-support tapes to functional tapes, using the native dense
relabelling. Then the new blank equivalence composes with
\texttt{KEndpoint215}. An endpoint for the new dense table also uses
FM215's supplied canonicalization map. These bridges must be explicit
if a new end-to-end Lean consistency endpoint is claimed. No new blank
bootstrap or calculus derivation is needed. Ordinary mathematical
composition with the established primary calibration already gives
the desired ZFC equivalence.

\subsection{Independent check of the existing CM215 route}

For comparison, \texttt{check\_cm215.py} in this packet independently checks
the existing six-row conditional-marker alternative against the exact
primary bytes. Its invariant is
\[
 q=52\Rightarrow u(h-1)=1,\quad
 (q=176\ \&\ u(h)=0)\Rightarrow u(h-1)=0,\quad
 q=156\Rightarrow u(h+1)=1.
\]
An old radius-two tape window determines the radius-one invariant after
one L/R transition. Enumerating all 32 such windows per working control
is therefore an exact finite proof-obligation check for arbitrary tapes.
Of 6,880 state/window combinations, 6,840 satisfy the invariant. All
preserve it and pass the whole-window block identity; 6,824 use one step
and 16 use three. The checker also verifies the supplied native/dense
bijection. It does not import the earlier symbolic runtime checker,
the earlier blank trace, or the source calculus. The ordinary locality
argument lifts these finite obligations to whole tapes because every
read and write lies inside the certified window and exterior cells are
untouched. The simpler FM215 invariant is the reason to prefer its new
formal proof as the next compact endpoint.


\section{A context certificate for the primary decoder}
\label{sec:decoder}

The next substantial shared-proof improvement is to replace repeated
prefix-level decoder checks by a certificate over shared source nodes
and the physical context that their implementation actually uses. It is
not enough to hash-cons the source tree alone: the same subtree can have
different physical entry controls, and some decoder paths inspect one or
two preceding bits. A proof must retain those distinctions.

The exact \path{sites215.json} input has 4,751 primitive leaves and 4,750
internal-prefix occurrences. Its maximally shared syntactic DAG has 202
nodes: 179 internal nodes and 23 distinct primitive labels including their
arguments. Pairing each internal node with its physical entry control
gives 180 roles. We then refine a role by the last at most two bits of the
sentinel-and-prefix string $1p$:
\[
 r=(v,q,\operatorname{suffix}_2(1p)).
\]
The resulting certificate has \textbf{234 roles}. At the root the sentinel
suffix is the one-bit string $1$ and the physical control is 84.

\begin{center}
\begin{tabular}{lr}
\toprule
Measured object & Count\\
\midrule
Expanded internal-prefix occurrences & 4,750\\
Shared internal source nodes & 179\\
Internal source-node/control pairs & 180\\
Pairs refined by one preceding bit & 196\\
Pairs refined by two preceding bits & 234\\
Internal branch rules in the two-bit certificate & 371\\
Primitive-entry branch rules & 97\\
Distinct primitive first-transition signatures & 68\\
\bottomrule
\end{tabular}
\end{center}

The larger number of roles after adding context is deliberate: source
sharing and physical correctness require different equivalence relations.
Two bits suffice for this certificate; no minimality assertion is made.
Refining by three, four or five bits would yield 296, 388 or 523 roles.
These counts are obtained from the actual literal table and source, not
from an estimate based on source-tree depth.

\subsection{The finite check and its whole-tape meaning}

For an internal branch, instantiate the role's suffix immediately to the
left of head position zero and put the branch bit at zero. Execute the
literal table until the head first reaches position one. The certificate
requires a positive number of steps, at most seven; every read and write
must stay inside the explicitly represented suffix and current cell;
every represented cell must be restored; and the final control must be
the target role's control. The target source node must be the appropriate
child and its suffix must be
\[
 \operatorname{suffix}_2(s\,b),
\]
where $s$ is the previous suffix and $b$ the branch bit.

The validator also checks the DAG's acyclic child order, expands it to the
exact pinned leaf list, checks the exact root role, rejects duplicate
roles, and requires every supplied role to be reachable from the root.
All two branches of every role are checked. For a primitive child the
rule binds its operation and argument to the literal first transition
from the role's control and branch bit.

\begin{proposition}[Internal traversal]
Suppose a full binary tape contains the sentinel and source address
prefix expected at an internal decoder role, with the head at the next
address bit. The internal rules of the supplied certificate take the
machine to the appropriate next role, preserve the entire tape, and
advance the head by one address cell in a positive finite number of
transitions. Iterating these rules reaches the correct primitive entry
for every source address.
\end{proposition}

\begin{proof}
For one internal branch, all inspected cells are in the explicit local
window. Agreement on that window makes the concrete full-tape run follow
the same transitions as the finite check. The full exterior is neither
read nor written. The check restores the window, so all tape cells are
preserved; its suffix equation establishes the next role's precondition.
Induct along the finite source address. The verified DAG child relation
ensures that each branch selects the corresponding source child. At a
leaf, the checked primitive-entry rule identifies its first transition.
No claim about execution beyond that first primitive transition is used.
\end{proof}

Across the expanded internal occurrences, there are 4,749 internal-child
edges: 4,355 take one transition, 303 take three, 90 take five, and one
takes seven. The maximum backward read is two cells. Thus the certificate
captures the real local decoder behavior, including its longer blocks.

The implementation is \path{checks/decoder_contexts.py}. The data and
receipt are in \path{evidence/decoder_context_certificate.json} and
\path{evidence/decoder_context_receipt.json}. This is an ordinary finite
certificate with a stated soundness argument. No Lean or Rocq theorem
for this new checker has yet been compiled.

\subsection{How to turn this into a smaller formal proof}

The proposed implementation sequence is concrete:
\begin{enumerate}
\item Define a finite source DAG and its expansion, and prove one generic
  equality between DAG traversal and tree traversal. Keep the present
  source semantics unchanged.
\item Define a role with source node, physical state and a suffix of at
  most two bits. Prove a generic local-frame lemma for functional tapes
  or the existing zipper interface.
\item Reflect each of the 371 internal and 97 primitive-entry rules into
  the kernel. Use the generic path-induction theorem to cover arbitrary
  address prefixes. This replaces repeated expansions of the same
  internal simulation facts by one reference to the role certificate.
\item Separately provide the existing quantified runtime contracts for
  the primitive entries, generalizing over the preserved address prefix
  and exterior as needed. The 68 first-transition signatures are useful
  grouping keys; they are not 68 already proved runtime contracts.
\item Compose these contracts at the shared site-coverage interface,
  retain the old endpoint as a regression target, and measure kernel term
  sizes and rebuild times before removing historical proof data.
\end{enumerate}

The present certificate establishes that substantial decoder repetition
can be shared soundly. It does not justify a numerical claim that the
whole correctness proof will shrink by $4750/234$, nor that the existing
guarded runtime obligations disappear. Its benefit must ultimately be
measured in the full build, not inferred from source counts alone.

\section{A practical proof architecture}
\label{sec:architecture}

The repository already shares much of its calculus and compiler proof.
The appropriate next architecture makes three interfaces explicit rather
than moving all certificates into a new monolithic development.

\paragraph{Machine equivalence for small mutations.}
Use \texttt{PositiveSimulation.halts\_iff} as a common interface for
transition-table edits. A candidate provides its invariant, state/configuration
relation, local positive blocks and terminal case. The endpoint then
reuses the parent's proof. FM215 is a compiled example. For a Rocq port,
the same induction on the remaining target halting time applies; it
requires neither an infinite trace nor an independent reverse simulation.

The direct route for incorporating FM215 into the existing Lean tree is
to identify the embedded native rows with the repository's K230 labels
(including $254\mapsto230$, $301\mapsto102$, $302\mapsto120$), prove the
functional-tape step relation equivalent to \texttt{KZip} or the core tape
semantics, and compose \texttt{KEndpoint215}. The existing
\texttt{KDenseBridge} already supplies the dense/zipper portion of the
parent interface. The new dense machine then uses its supplied finite
renaming. This adapter is the remaining mechanical integration task; it
does not require redoing the calculus or the 95,714-step initialization.

\paragraph{Finite structural certificates for compiler output.}
Use an indexed source DAG and context roles for the decoding skeleton;
retain quantified contracts for register primitives. Prove the checker
sound once and evaluate only small, machine-specific propositions. Bind
each theorem to a literal table and source representation. A file hash is
an excellent audit key, but it is not a substitute for the literal data
or an equality proved in the kernel.

\paragraph{A stable calibration interface.}
Keep the numerical-calculus/ZFC equivalence behind its named theorem.
Expose explicit hypotheses for numerical guards and encoding side
conditions. Do not repeat the already completed shared AX interpreter,
Lean mirror-certificate removal, or guarded dominated-rule refactor.
The remaining Rocq mirror theorem is the smallest immediate source-level
change; the decoder certificate is the larger structural project.
\cite{phase3,source215,offsetbisim,leansim}

For one or two alternatives, direct simulation is likely to be much less
work than proving a general optimizing compiler. If many future machines
will use arbitrary reallocation and reencoding, a verified certificate
checker becomes more attractive: its initial proof cost is amortized over
new tables. Candidate selection should therefore compare three quantities:
physical state count, number of distinct proof contexts, and strength of
the required tape invariant. A source tree with fewer nodes or shorter
addresses does not automatically win on the latter two.

% Include-only mathematical note. Requires amsmath, amssymb, amsthm.
\section{An obstruction to the large-matrix logical recovery}
\label{sec:tail-quantifier-obstruction}

The existing large-matrix experiment removes the final comparison gates by
emitting the seed bundle only at parameters $c>a+b$. Its report explicitly
leaves recovery of the missing atomic-vacuity interface open. The following
result closes one proposed route: that interface cannot be recovered from
the retained logical rows alone. It does not decide whether the seed theory
itself supplies a different consistency reduction.

Write $R(a,b)$ and $E(a,b)$ for membership and equality atoms, and
$Q_i A$ for the named universal-quantifier constructor. Let
$\mathcal L_{>}$ contain truth, modus ponens, theorem generalization, the
unrestricted classical implicational basis, the rows $\rho_1,\ldots,\rho_6$
of \texttt{Code254.lean}, and the merged equality row
\[
 E(c,a)\to\bigl(E(c,b)\to Q_{r(a,b)}E(a,b)\bigr),
 \qquad r(a,b)=4P(a,b)+1.
\]
Its membership-vacuity axioms are only
\[
 R(a,b)\to Q_cR(a,b) \qquad(c>a+b).
\]
Allowing the unrestricted classical basis makes the obstruction stronger:
the accepted dominated $\rho_8$ instances are a subfamily of it.
No seed-K bundles or other nonlogical axioms are included in
$\mathcal L_{>}$.

\paragraph{Theorem.}
The formula $R(2,3)\to Q_1R(2,3)$ is not derivable in
$\mathcal L_{>}$. In particular, the retained logical subsystem does not
recover atomic vacuity at every distinct binder.

\paragraph{Proof.}
For any nonempty set $D$ and irreflexive binary relation $R_D$, take worlds
to be all sequences $e:\mathbb N\to D$. Interpret atoms by
\[
 e\models R(a,b)\iff R_D(e(a),e(b)),\qquad
 e\models E(a,b)\iff e(a)=e(b).
\]
Interpret implication and falsity classically, and replace ordinary
single-variable quantification by \emph{tail quantification}:
\[
 e\models Q_i A
 \iff
 \forall f\bigl[(\forall k<i,\ f(k)=e(k))\Longrightarrow f\models A\bigr].
\]
The relations $e\sim_i f\iff\forall k<i,\ e(k)=f(k)$ are equivalence
relations. They are nested, and
\[
 \sim_i\circ\sim_j=\sim_{\min(i,j)},\qquad
 Q_iQ_jA\ \equiv\ Q_{\min(i,j)}A.
\]
This identity is an identity of the semantics for every formula $A$;
it is not a new inference rule.

All retained logical rows are valid. Normality of universal
quantification proves $\rho_2$, and the displayed identity proves
$\rho_4$. The row $\rho_6$, $(\exists_iQ_iA)\to A$, is valid because
if an $i$-accessible world satisfies $Q_iA$, then the original world is
among its $i$-accessible worlds. The row $\rho_1$, $\exists_i E(i,j)$,
is valid: if $j<i$, assign coordinate $i$ the fixed value $e(j)$;
otherwise assign coordinates $i,j$ a common value of $D$. Both changes
preserve the prefix below $i$. Equality substitution proves $\rho_3$
and $\rho_5$. For the merged equality row, $r(a,b)>a,b$, so
$r(a,b)$-accessible worlds preserve coordinates $a,b$; the two equality
premises imply their equality there as well. The classical
implicational axioms are tautologies. Global validity is closed under
modus ponens and under every $Q_i$, as required by theorem
generalization.

If $c>a+b$, then $a,b<c$, so $c$-accessible worlds preserve both
arguments of $R(a,b)$. Thus all permitted atomic-vacuity rows are valid.
The designated numerical falsity $R(0,0)$ is false everywhere, since
$R_D$ is irreflexive. The numerical truth alias $I(0,0)=0$ is harmless:
both sides evaluate to true. Equivalently, compose this interpretation
with the repository's recursive numerical decoder; its implication and
quantifier equations remain valid with these semantics.

Finally choose $D=\{0,1\}$ and $R_D(x,y)\iff x<y$. Let $e(2)=0$
and $e(3)=1$. Then $e\models R(2,3)$. There is a world $f\sim_1 e$
with $f(2)=f(3)=0$, so $f\not\models R(2,3)$ and
$e\not\models Q_1R(2,3)$. Hence the proposed implication is not valid
in a semantics validating every axiom and inference rule of
$\mathcal L_{>}$, and therefore is not derivable.\hfill$\square$

\paragraph{What this resolves and what it leaves open.}
Ordinary first-order matrix padding, or an unproved appeal to arbitrary
variable renaming, cannot complete the missing logical bridge. A proof
that transforms an allowed high-binder vacuity instance into the displayed
low-binder instance must use additional nonlogical assumptions or a rule
not present in $\mathcal L_{>}$. This result does not assert that the
full large-matrix calculus has this countermodel: its seed bundles have
not been shown valid in tail semantics. It neither proves nor disproves
that full calculus's equiconsistency with ZFC, and gives no state-count
lower bound.

\paragraph{Executable finite audit.}
\texttt{check\_tail\_quantifiers.py} checks the counterexample on the
16 assignments of coordinates $0,1,2,3$, all indicated atomic rows with
parameters in that range, and the universal algebraic laws used for
quantifiers. It checks normality from the finite accessibility relations,
not by sampling formulas, and checks $\rho_6$ and commutation for all
65,536 propositions on those 16 worlds. The finite audit is an independent
sanity check of the displayed example and semantics.

\paragraph{Compiled Lean interface.}
\texttt{TailQuantifiers.lean} additionally formalizes the all-natural
tail semantics, the validity of each displayed logical row, soundness of
an inductive proof system, and the explicit non-derivability conclusion.
It imports only \texttt{Init} and was compiled with Lean 4.32.0. Its
logical interface is stronger than the one needed here: it admits every
membership-vacuity binder larger than both atom indices, and every merged
equality helper binder larger than both equality indices. The theorem
\texttt{not\_valid\_missing\_vacuity} has no axioms; the soundness and
non-derivability theorems use only \texttt{propext},
\texttt{Classical.choice}, and \texttt{Quot.sound}. There are no
\texttt{sorry} or \texttt{admit} declarations. The file treats named
logical formulas as predicates on worlds; the extension through the
repository's numerical decoder is the ordinary constructor argument
above, not a newly compiled numeric-calculus bridge. Neither the Lean
file nor the finite audit includes the seed-K theory bundles.

\paragraph{Pinned repository sources.}
This note concerns the pinned ProveIt revision, in particular its
large-matrix formalism supplement\cite{largematrix},
\path{BusyBeaver/ZFC/Code254.lean},
\path{BusyBeaver/ZFC/KCode222.lean}, and
\path{BusyBeaver/ZFC/KThmK.lean}. The logical schemas above are their
literal named-formula interpretations. This is a new obstruction proof;
the large-matrix source experiment and the numerical calculus are inherited
work.


\section{What the search toward 214 establishes}
\label{sec:search}

The search uses the primary K215, thirteen pinned alternatives and FM215.
It begins with a sound local invariant: for each state, collect reachable
seven-cell windows centered at the head. Initially state zero admits all
128 windows, representing arbitrary initial tapes. Under each literal
transition, shift the window and admit both values of the newly exposed
exterior bit. Taking the least fixed point gives an overapproximation of
the concrete reachable windows. This abstraction is inductive but loses
long-range tape correlations.

The broadest completed family examines pairs of noninitial controls with
identical other-symbol transitions and different selected-symbol
transitions. For that selected symbol, the two sets of admitted seven-bit
windows must be nonempty and disjoint. It deletes one of the pair,
redirects every incoming edge to the survivor, and tries each write,
direction and surviving working-state target for the differing row;
halt is excluded as a replacement target. Both choices of survivor are
tested. The symbolic comparison is required to match the parent's
mapped endpoint, head and all inspected tape cells for every admitted
seven-cell window, branching on previously unknown exterior cells.

\begin{center}
\begin{tabular}{lrr}
\toprule
Search family & Cases & Successful candidates\\
\midrule
Full-window differing-row replacements & 1,761,648 & 0\\
Adjacent-marker subfamily (included above) & 710,480 & 0\\
Separate two-helper-row resynthesis & 88,395 & 0\\
\bottomrule
\end{tabular}
\end{center}

The broad family has 1,029 eligible control pairs. Its replacement count
is $1029\cdot2\cdot(2\cdot2\cdot214)=1{,}761{,}648$ across the 15
tables. The symbolic exploration limits are 32 transitions, 13 inspected
cells and 200 symbolic nodes, with the initial seven cells and at most six
new exterior cells. The adjacent-marker subfamily uses 415 pairs and an
eight-cell inspection limit. These are nested counts, not independent
searches to add together. An earlier 190,032-case read-zero/right-moving
search is also a subset and is not counted again.

The separate helper search checks $415\cdot213=88{,}395$ pair/helper
choices. It allows both rows of one helper to change, imposes the exact
three-step selector effects, and preserves old helper uses through a
two-step equality on every symbol admitted by the window invariant.
Those restrictions define its scope. An exploration-limit exit is
unresolved, not a counterexample. A mismatch on an abstract window can
also be caused by the overapproximation; it does not by itself reject the
candidate's actual blank-tape behavior.

The retained scripts and per-table receipts make the family and counts
inspectable. No search ranges over all 214-state tables. In particular,
source changes, global role reallocations, richer tape invariants and
longer simulation blocks remain outside these families. The result is no
verified candidate, not evidence of a 215-state lower bound.

One promising fold has a more informative failure than a bound exit. In
the commands120 table, merge controls 78 and 210, redirect entries to 78,
and replace $78/0$ by $0R201$. On source leaf 25, address
$1101001000$ and command CD3, take all nine registers zero, one zero spare
and gap zero. The parent returns to the correct successor address after
610 transitions; this 214-state trial returns after 594 transitions with
the one at address-cell position 8 lost and a one spuriously added at
position 11. The supplied independent checker verifies the valid input
frame, exact source leaf, both full output tapes and the literal trial.
This refutes the proposed uniform primitive-runtime contract. It does not
by itself decide the trial's blank-run equivalence. Its table is explicitly
retained as rejected research under \path{machine/research/}.

\subsection{Where to search next}

FM215 isolates the unresolved resource question precisely: its logical
role merger consumes the freed control as the predecessor selector. To
save a physical state, that selector behavior must be implemented by an
already live helper without breaking the helper's old uses, or several
roles must be reallocated together.

The next useful search is therefore a multi-row synthesis problem. Fix a
small neighborhood, allow two or more helpers to change, give each its
reachable tape contexts, and ask for positive local simulation blocks.
Use symbolic counterexamples to refine the contexts. For every apparent
winner, produce a certificate accepted by the same generic simulation
interface used here. An SMT or SAT solution is a proposal generator; the
small kernel theorem should remain the acceptance criterion.

A second route changes the source or proof enumeration. It can make a
larger physical difference, but it must carry a consistency-calibration
argument before its state count is meaningful. Section~\ref{sec:tail-quantifier-obstruction}
shows why one tempting removal of numerical comparison gates cannot be
justified by logic-only vacuity recovery. It leaves theory-specific
translations open. Startup shortening alone is also insufficient: changing
the initial entry without preserving its complete calibrated tape creates
a new global invariant obligation.

\section{Novelty, reproducibility and incorporation}
\label{sec:reproduction}

The four-row construction was compared with every \texttt{.tm} path in
the pinned ZFC tree: 622 paths, with no missing or unparsable table. Of
these, 321 are complete graph-reachable 215-state tables, representing
84 distinct rooted transition graphs. FM215 matches none, even allowing
global left/right reflection. \path{machine/novelty_current_tree.json}
records each current Git blob and SHA-256. This is novelty relative to
that exact corpus; no claim about all unpublished machines or priority
outside the repository is made.

The dense FM215 table has SHA-256
\begin{verbatim}
1b602d3b37887060ce2d992cf3036f540e8ccff5d37ab06c80f19f1edd38cbc6
\end{verbatim}
Its 95,840-step blank prefix contains 63 stretched parent steps, giving
$95{,}714+2\cdot63=95{,}840$. At that boundary the whole tape, not just a
bounded observation window, equals the parent's tape: dispatcher 84,
head $-8$, 95 one-cells and zeros elsewhere. The inherited interpretation
is registers $[4,0,0,0,0,0,0,0,0]$, gap one, and forty unit spares.
This agreement is checked by the independent sparse-tape simulator.

\subsection{Files and commands}

\begin{center}
\begin{tabular}{>{\raggedright\arraybackslash}p{.24\textwidth}>{\raggedright\arraybackslash}p{.68\textwidth}}
\toprule
Location & Contents\\
\midrule
\path{report.tex}, \path{sections/} & This report's complete editable source.\\
\path{report.pdf} & Compiled and visually inspected report.\\
\path{formal/} & Three FM215 Lean modules, exact native tables, compiler
logs, receipt, binding checker, and independent CM215 comparison.\\
\path{machine/} & Native/dense FM215, renaming, builder, startup/local
checks, pinned comparison tables, novelty and search scripts/receipts.\\
\path{checks/}, \path{evidence/} & Decoder context certificate and audit;
bundle verification entry point and receipts.\\
\path{research/tail-quantifiers/} & Lean obstruction proof, finite audit,
compiler log and exact source/compiler receipt.\\
\path{inputs/}, \path{provenance/} & Exact source inputs, inherited
licenses/notices, and pinned upstream evidence references.\\
\bottomrule
\end{tabular}
\end{center}

From the packet root, the following commands require no ProveIt checkout
and no network access after the tools are installed:
\begin{verbatim}
python3 -B verify_all.py
python3 formal/check_fm215.py --lean /absolute/path/to/lean
lean research/tail-quantifiers/TailQuantifiers.lean
pdflatex -interaction=nonstopmode -halt-on-error report.tex
pdflatex -interaction=nonstopmode -halt-on-error report.tex
\end{verbatim}
The default verifier uses Python's standard library and rechecks the
construction, exact literal bindings, local/startup simulations, decoder
certificate, CM215 comparison, finite countermodel and receipt source
hashes. It does not silently describe saved compiler logs as a fresh
kernel run. The explicit Lean commands replay the kernel proofs and
require Lean 4.32.0; no Mathlib project is required. The FM215 driver sets
its local import path. Generated compiled objects are intentionally not
part of the portable source packet.

The full novelty audit additionally needs the pinned ProveIt Git tree.
The two search scripts use standard Node.js and are optional, longer
reproductions; their precise commands are in the top-level README. The
archive manifest binds the shipped files. Rebuilding logs or the PDF can
change their hashes, so compare the manifest before regeneration.

\subsection{Recommended repository integration order}

First add FM215 and its standalone equivalence as a named alternative,
with the explicit ordinary composition and the remaining Lean adapter
identified. Then complete that adapter and, if desired, port the small
positive-simulation proof to Rocq. Next port Rocq's mirror-certificate
symmetry. Finally integrate the context certificate behind the existing
site-coverage interface and benchmark the full dependency closure.
This order yields a small reviewable theorem at each step and preserves
the current certified parent as a stable comparison point.

The parent construction, numerical calculus, ZFC calibration, compiler
ancestry and previously accepted alternatives remain inherited work.
The packet's new contributions are the FM215 four-row transformation and
checked equivalence, the decoder context certificate, the named logical
obstruction and its proof, and the bounded search/audit work. Original
notices and licenses are retained. The report does not claim a fresh full
rebuild of the repository's existing Lean or Rocq development.

\begin{thebibliography}{10}
\bibitem{verification}
ProveIt, \repo{Computability/BusyBeaver/Research/ZFC/k225-formal/verification/k215/README.md}{K215 formal verification record}, pinned revision above.
\bibitem{phase3}
ProveIt, \repo{Computability/BusyBeaver/Research/ZFC/k225-formal/verification/k215/simplification/phase3/README.md}{Phase-three proof-simplification record}.
\bibitem{primary}
ProveIt, \repo{Computability/BusyBeaver/Research/ZFC/candidate-k215/README.md}{Primary K215 candidate and provenance}.
\bibitem{endpoint}
ProveIt, \repo{Computability/BusyBeaver/Lean/BusyBeaver/ZFC/KEndpoint215.lean}{\texttt{KEndpoint215.lean}}.
\bibitem{leansim}
ProveIt, \repo{Computability/BusyBeaver/Lean/BusyBeaver/ZFC/KSim.lean}{\texttt{KSim.lean}} and
\repo{Computability/BusyBeaver/Lean/BusyBeaver/ZFC/RMOffsetBisim.lean}{\texttt{RMOffsetBisim.lean}}.
\bibitem{source215}
ProveIt, \repo{Computability/BusyBeaver/Research/ZFC/k225-formal/rocq/rocq-extension/coq/215/Source215.v}{Rocq \texttt{Source215.v}}.
\bibitem{offsetbisim}
ProveIt, \repo{Computability/BusyBeaver/Research/ZFC/k225-formal/rocq/rocq-extension/coq/215/OffsetBisim230.v}{Rocq \texttt{OffsetBisim230.v}}.
\bibitem{cm}
ProveIt, \repo{Computability/BusyBeaver/Research/ZFC/candidate-k215/alternatives/conditional-marker/README.md}{Conditional-marker alternative}.
\bibitem{largematrix}
ProveIt, \repo{Computability/BusyBeaver/Research/ZFC/candidate-k215/alternatives/rebalanced/research/formalism/formalism_supplement.tex}{Large-matrix formalism experiment}.
\bibitem{simplification}
ProveIt, \repo{Computability/BusyBeaver/Research/ZFC/k225-formal/verification/k215/simplification/README.md}{Proof-simplification checkpoints and source-byte comparisons}.
\end{thebibliography}
\end{document}
